\cotefile{161-2}
\scanpage{66}
Existence de $\varprojlim$ ou $\varinjlim$ dans $T(C)\inserted{\simeq\underline{\mathrm{Hom}}_{\Delta}(R,C)}$, pour $T$ \unsure{2-représentable|2-representable}.

Soit $\delta$ un type de diagramme, tel que

a) Les $\varprojlim$ (disons) de type $\delta$ existent dans $\mathbf C$

b) les $\varprojlim$ de type $\delta$ commutent, \deleted{pour tt $C'\in\mathrm{Ob}\,\Delta$,} \inserted{dans $C$} aux $\unsure{\lim|\varinjlim}$ de type $d$ (i.e. aux $\varinjlim$ de type $d''$)

Alors les $\varprojlim$ de type $\delta$ existent dans $T(C)$, et se calculent dans $\underline{\mathrm{Hom}}_{\Delta}(R,C)$ ``\uline{argument} par \uline{argument}'' (en d'autres termes, les foncteurs $r_{C}:T(C)\longrightarrow C$ associés aux \inserted{$T$-}constructions commutent aux $\varprojlim$ de type $\delta$)

De plus, pour un hom de théories \unsure{repr.} $T\longrightarrow T'$, $T(C)\longrightarrow T'(C)$ commute aux $\varprojlim$ de type $d$. Id pour $C\longrightarrow C'$, $C'$ satisfaisant les $\unsure{\text{mêmes}|\leftrightarrows}$ hyp. que $C$, $T(C)\longrightarrow T(C')$ commute aux $\varprojlim$ de type $d$.

On a, pour $R,\unsure{\mathscr C|\mathbf C}\in\mathrm{Ob}\,\Delta$, \deleted{\unsure{$\mathrm{Ob}\,\mathcal C$}} et $\unsure{\mathbf J|\mathbf I}\in\mathrm{Ob}\,\mathrm{Cat}$
\[
\underline{\mathrm{Hom}}_{\mathbf{Cat}}(\mathbf J,\mathscr C)\in\Delta
\qquad\text{(au moins si $\mathbf J$ petite)}
\]
\[
\underline{\mathrm{Hom}}_{\Delta}(R,\underline{\mathrm{Hom}}_{\mathrm{Cat}}(\mathbf J,\mathscr C))
\simeq\underline{\mathrm{Hom}}_{\mathrm{Cat}}(\mathbf J,\underline{\mathrm{Hom}}_{\Delta}(R,\mathscr C))
\]
i.e.
\[
\boxed{T_R(\underline{\mathrm{Hom}}_{\mathrm{Cat}}(\mathbf J,\mathscr C))\simeq
\underline{\mathrm{Hom}}_{\mathrm{Cat}}(\mathbf J,T_R(C))}
\]

Faisons p. ex. $\mathscr C=(\mathrm{Ens})$ (supposant $(\mathrm{Ens})\in\mathrm{Ob}\,\Delta$ !), $\mathbf J$ remplacé par $\mathbf J^{\circ}$
\[
T_{\unsure{\mathscr R|R}}(\widehat{\mathbf J})\overset{\sim}{\simeq}\underline{\mathrm{Hom}}(J^{\circ},T_R(\mathrm{Ens}))
\]

Si $C\longrightarrow C'$ dans $\Delta$ est pl. fidèle (resp. fid) alors $T(C)\longrightarrow T(C')$ l'est aussi (pour $T$ 2-repr.). \deleted{et} Si $T\longrightarrow T'$ est associé à $R'\longrightarrow R^{\unsure{\prime\prime|\prime}}$ ``\inserted{$\Delta$-}générateur'' (tel que, donc, tte sous-catégorie pleine $\in\Delta$ de $R'$ par \unsure{lequel|laquel|laquelle} ça se factorise soit \deleted{\unsure{dans}} $\simeq R$] alors le diagramme
\diagram{\begin{tikzcd}[ampersand replacement=\&]
T(C) \arrow[r,hook] \arrow[d] \& T(C') \arrow[d] \\
T'(C) \arrow[r,hook] \& T'(C')
\end{tikzcd}}
est 2-cartésien
% archive sheet number: 65
